%% 第二章

\chapter{相关技术与理论基础}
\chaptermark{相关技术与理论基础}

本章为后续方法描述提供术语、公式与评价指标基线。本章不引入新的实验结论，仅围绕已有文献组织技术铺垫。

\section{B-rep 与 NURBS 几何表示}

\subsection{B-rep 拓扑层次}

B-rep 实体由 Solid、Shell、Face、Loop、Edge、Vertex 六层拓扑组成。Shell 提供方向（朝外或朝内），Face 通过其外环和若干内环对参数曲面进行裁剪，Wire 是 Edge 的有序集合，Edge 由几何曲线 + 起止 Vertex 描述。这种层级表示既保留了局部几何（曲面与曲线方程），又保留了全局拓扑（面—边邻接、边—顶点共享）\cite{mantyla_1988_solid_modeling,hoffmann_1989_geometric_solid_modeling,piegl_tiller_1997_nurbs}。这种“几何与拓扑分离”是 B-rep 区别于体素或网格的根本特征，本文方法也正是利用这一特征——分别为面与边设计采样规则，再用面索引把几何与拓扑串成统一序列。

在几何内核中，B-rep 的合法性并不只取决于每一张 Face 是否能被单独画出，还取决于不同层级之间的引用关系是否一致。最常被忽视的一类约束是边的二维与三维表示必须匹配：一条 Edge 通常同时由三维空间曲线和它在相邻 Face 参数域上的二维 p-curve 描述，前者负责给出边在世界坐标系中的位置，后者负责把这条边落到具体面上的 UV 区域上，当两者通过容差检查存在偏差时，实体在视觉上接近闭合却会被内核判为无效。另一类常见约束是 Loop 方向与 Shell 法向的一致性：Loop 方向决定参数域上的裁剪内外，Shell 中各 Face 的法向决定 Solid 的内外体积，二者一旦反转，BRepCheck 可能报告曲线不在曲面上、取向错误或壳体不闭合等问题。换言之，直接生成 B-rep 必须同时关注三个层次的一致性：“面是什么形状”、“边如何连接面”、“整体是否能围成封闭体”，缺一不可。本文逆向重建管线之所以把 NURBS 拟合、面索引解析、拓扑缝合与形状修复串成一条线，也是为了让这三层一致性的检查位置足够清楚。

本文序列化设计将面段放在边段之前，也与这种层级关系有关。面段首先提供每个面的几何样本与面索引，使模型建立一组局部曲面片的先验，边段随后通过两个邻接面索引把这些面联系起来，再附带边曲线采样点约束相邻面的交界位置。虽然该表示不是完整 CAD 内核数据结构的逐字段复制，但它保留了 B-rep 生成中最关键的两个信息源：NURBS 几何采样和面—边邻接关系。后续重建时，OpenCASCADE 根据这些信息重新拟合曲面、曲线并执行缝合，从而把 token 序列还原为可以被几何内核检查的实体对象。

\subsection{NURBS 曲线与曲面定义}

NURBS 曲线和曲面由控制点、节点向量、权重三组参数刻画。B 样条基函数通过 Cox--de Boor 递推定义：当 $p=0$ 时

\begin{equation}
    N_{i,0}(u) = \begin{cases} 1, & u_i \le u < u_{i+1} \\ 0, & \text{otherwise} \end{cases}
    \label{eq:bspline-base-0}
\end{equation}

当 $p \ge 1$ 时

\begin{equation}
    N_{i,p}(u) = \frac{u - u_i}{u_{i+p} - u_i} N_{i,p-1}(u) + \frac{u_{i+p+1} - u}{u_{i+p+1} - u_{i+1}} N_{i+1,p-1}(u).
    \label{eq:bspline-base-p}
\end{equation}

NURBS 曲面 $S(u, v)$ 由两组方向上的基函数与控制网格 $\mathbf{P}_{ij}$、权重 $w_{ij}$ 共同决定（见\autoref{eq:nurbs-surface}）。在本文实验中，预处理阶段对每个面在 UV 参数域上采样若干三维点构成面段，对每条边在弧长方向上采样若干三维点构成边段，具体的采样数量与分块组织方式见第~3 章。

NURBS 的一个关键工程性质是局部支撑：节点向量上某个区间的基函数仅在该区间附近不为零，因此单个控制点只影响其邻域的曲线或曲面形状，而不像全局多项式那样牵动整段曲面。该性质对工业 CAD 建模和本文生成任务都有直接意义。对建模而言，设计者可以独立调整局部倒角、孔洞、曲面过渡而不破坏其他区域。对生成任务而言，它意味着曲面几何可以通过有限数量的采样点近似刻画。只要采样密度足以反映主要曲率变化，后续 B 样条曲面拟合阶段就有机会恢复出可用的曲面片。基于此，本文为每个面取 16 个三维采样点、为每条边取 4 个三维采样点。这一密度并不是由理论给出的最优值，而是在序列长度上限、几何细节保留能力与本科毕设可用算力之间的工程折中：若把每个面提高到 25 或 36 个采样点，第~5 章序列长度直方图上的尾部样本会更频繁触及 4\,096 token 的长度上限。若降到 9 个点，平面以外的曲面片就难以保留主曲率。


\begin{figure}[!htbp]
    \centering
    \includegraphics[width=0.82\textwidth]{插图/fig2-1_nurbs_sampling.pdf}
    \caption{NURBS 曲面与曲线采样示意图。左侧展示曲面 $S(u,v)$ 在 UV 参数域上的规则采样过程，每个面采样 16 个三维点，并按规范化方向展开为面段的坐标部分。右侧展示曲线 $C(u)$ 沿弧长方向的边采样过程，每条边采样 4 个三维点，并通过端点坐标统一边方向。采样得到的三维坐标随后被标量量化为 4000 个等级，并映射到大小为 4053 的数据侧统一词表中。}
    \label{fig:nurbs-sampling}
\end{figure}

\subsection{CAD 文件交换格式}

STEP（ISO 10303）是工业 CAD 标准交换格式，可直接存储 B-rep 拓扑、NURBS 曲面与曲线参数，其产品数据表达与交换体系在国内标准中也有对应规范\cite{gbt16656_1_2008,gbt16656_11_2010,gbt16656_21_2008}。STL 仅存三角网格，常用于 3D 打印或可视化。本文最终输出 STEP 文件以保留几何精度，STL 仅作为可视化辅助。STEP 的读写、有效性检查与导出均通过 OpenCASCADE 完成，相关流程在第~4 章说明。

从论文评测角度看，选择 STEP 而非 STL 作为主输出格式也会提高任务难度。STL 文件只要求三角面片可以被写出和渲染，很多局部自交、法向不一致或边界缝隙问题在简单可视化时不一定暴露，STEP 文件则需要几何内核能够解析其拓扑结构，并进一步接受 BRepCheck 等规则检查。本文把 STL 放在可视化辅助位置，是为了便于人工观察生成零件外观，但所有有效率数字均以 STEP / B-rep 校验为准，避免用“看起来像”替代“内核认为合法”。

\section{CAD 数据集与生成方法}

\subsection{常用 CAD 数据集}

CAD 领域已有多个公开数据集为生成式研究提供了素材，但它们的来源、规模和数据形式差异显著。ABC 数据集由百万级机械零件 CAD 组成，覆盖了从家用器件到工业组件等广泛形态，几何复杂度与拓扑形态的覆盖面较广\cite{koch_2019_abc}，但内部样本质量参差不齐，常含有读取失败或拓扑退化的实体，直接使用需要较重的清洗。DeepCAD\cite{wu_2021_deepcad} 在 ABC 基础上经过过滤并把原始 STEP 还原为草图—拉伸命令序列，最终提供约 18 万个可执行命令脚本，是当前 CAD 生成研究中使用最广的基线数据集之一。Fusion 360 Gallery 与 SketchGraphs 分别从人类 CAD 设计的命令历史和约束图角度提供补充数据来源\cite{willis_2021_fusion360,seff_2020_sketchgraphs}，更接近真实设计行为但样本规模相对较小。在以 B-rep 为输入的判别 / 表征学习方向，UV-Net\cite{Jayaraman_2021_CVPR} 与 BRepNet\cite{lambourne_2021_brepnet} 等工作进一步说明，面、边及其拓扑邻接关系不仅是 CAD 内核内部结构，也可以作为神经网络的输入结构，由 GNN 在 B-rep 上做消息传递。本文在主实验中采用 DeepCAD 作为唯一数据集（详 §5.1），主要原因是它的规模、清洗程度与可重复性更适合本科毕设阶段建立稳定基线，且能保证后续与其他自回归基线在数据口径上具备初步可比性。


此外，本文使用的是由 DeepCAD 零件转换得到的 B-rep / STEP 形式，并在预处理阶段重新抽取 NURBS 采样点。这意味着训练目标不再是原始 DeepCAD 论文中的草图—拉伸命令序列，而是与本文第~3 章定义一致的点级 token 序列。因此，虽然数据来源与 DeepCAD 相关，本文模型的输出空间、损失函数和评测方式都与 DeepCAD 原始任务不同。第~5 章中的有效率结果只能说明本文表示和工程管线在该子集上可行，不能直接与 DeepCAD 的命令序列指标做逐项等价比较。

\subsection{B-rep 生成方法分类}

如第~1 章所述，现有 B-rep 相关生成方法大体可以归纳为以下三类方法，它们各自的中间表征与建模方式不同。

(1) CAD 操作序列建模。以 DeepCAD\cite{wu_2021_deepcad} 为代表，将设计过程还原为草图、拉伸、布尔等参数化 CAD 操作序列，并用 Transformer 自回归生成操作命令。其生成对象是“操作脚本”而非 B-rep 几何本身，需经过 CAD 内核执行才能得到 Solid。

(2) 草图—拉伸解耦 + 离散码本。以 SkexGen 为代表，将草图与拉伸表征分别投影到多个相互解耦的离散码本，再在每组码本上独立训练自回归生成器，最后通过联合采样把不同码本上的索引拼成完整的 sketch-extrude 序列。这种解耦让“画什么草图”和“如何把草图拉伸成特征”两个子分布相互松绑，提升了采样多样性与组合自由度。但建模目标依然停留在 sketch-extrude 操作脚本层面，并未触及 B-rep 几何采样点，同时几何精度由码本规模决定，进一步压缩几何细节需要扩大码本，反过来又会让稀疏码本难以训练。

(3) 结构化潜空间扩散。以 BrepGen\cite{xu_2024_brepgen} 为代表，直接对 B-rep 的面、边、顶点编码为结构化潜变量，并用扩散模型在潜空间中联合采样几何与拓扑。该路线不需要 CAD 操作字典，几何与拓扑可一并生成，但需要预训练几何编码器，且潜变量分布的精确拟合对生成质量有较大影响。

本文不复现上述任何 baseline，而是仅引用文献报告值作为口径对照。本文方法与上述三类方法在表征上的差异是：将量化下移到坐标级，不引入显式 VQ-VAE 码本，也不通过扩散在潜空间采样，而是直接在“点级 token + 自回归 Transformer”上完成生成。

这种表征选择带来一个明显特点：序列长度会变长，但几何量化误差更容易解释。VQ-VAE 码本方法把局部几何压缩到离散索引，模型生成的是码本编号，最终几何质量受编码器、码本容量和解码器共同影响，本文把坐标分量直接量化为 4000 个等级，单个 token 的含义更接近几何空间中的一个标量位置。这样做牺牲了压缩率，却减少了“码本索引与真实几何之间”这一层不可见误差源。对于以 BRepCheck 通过率为主目标的原型系统而言，这种可解释的误差来源更便于后续定位问题。

\subsection{多边形网格自回归方法}

PolyGen\cite{nash_2020_polygen} 是与 B-rep 路线相邻的另一类直接序列方法，它把三角网格的顶点和面索引分别视为可被自回归 Transformer 生成的离散序列对象。该工作的工程意义在于实证了“多 token 自回归生成可以稳定建模结构化几何”，包括变长序列、稀疏依赖与离散索引等问题。本文沿用这一基本思路并扩展到 B-rep：与 PolyGen 用两个 Transformer 分别生成顶点和面索引不同，本文把面 / 边的几何采样点与拓扑面索引放进同一统一词表、由单一 decoder-only Transformer 顺序生成，token 类别由词表段落位置自然区分。这样既避免了多模型协同的复杂度，也使面索引与坐标 token 之间的相互约束可以由同一注意力机制直接建模。

\section{自回归序列建模}

\subsection{Decoder-only Transformer 架构}

Transformer\cite{vaswani_2017_attention} 提出多头注意力机制，通过查询、键、值的点积匹配实现全局上下文聚合。decoder-only 变体（如 GPT、LLaMA 系列）使用因果掩码确保第 $t$ 个位置只能看到 $t$ 之前的 token，从而支持自回归生成。这一架构已经成为大规模语言建模的事实标准，也被引入到图像、语音、生物序列等领域。

\begin{figure}[H]
    \centering
    \includegraphics[width=0.74\textwidth]{插图/fig2-2_transformer_decoder.pdf}
    \caption{Decoder-only Transformer 结构示意图。输入 token 经嵌入与位置编码后进入重复堆叠的 Decoder block，每个 block 在因果掩码约束下包含 masked multi-head self-attention、前馈网络、归一化与残差连接，最终隐藏状态经线性预测头和 Softmax 得到下一个 token 的概率分布。}
    \label{fig:transformer-decoder}
\end{figure}

\subsection{自回归训练损失}

自回归模型最大化训练序列的对数似然，等价于最小化每个位置的交叉熵：

\begin{equation}
    \mathcal{L} = -\sum_{t=1}^{T} \log p(y_t \mid y_{<t}, \theta).
    \label{eq:ar-loss}
\end{equation}

训练时采用 Teacher Forcing：将真实 token 序列 $y_{<t}$ 直接作为输入，单步并行预测所有位置，推理时则按一个时间步采样、回填、再推进的方式串行展开。损失计算时跳过 PAD 填充位置，使其不参与梯度更新。

对本文任务而言，Teacher Forcing 的优点是训练效率高：所有位置的 next-token 预测都可以在一次前向中并行计算，损失也可以平均到所有非 PAD 位置。代价是训练阶段与推理阶段存在暴露偏差：训练时模型每一步都看到真实历史 token，推理时却只能看到自己上一步采样得到的 token。一旦早期一步出现面索引或 SEP 位置错误，错误就会成为推理上下文的一部分并被持续放大，使序列偏离合法骨架。本文从两条路径缓解这一问题。其一是在数据侧使用固定长度的面段和边段骨架：每个面段恒长 49，每个边段恒长 14，并以 START / SEP / END 严格分段，让模型在训练分布中反复看到稳定模式，提高其对位置类型的识别能力。其二是在推理与重建阶段设置最大序列长度、END 终止条件以及解析失败的明确记录，避免重建管线把已经偏离骨架的序列强行解释为某种实体。这部分判别由 §1.2.4 给出的三态口径承担。更强的语法约束解码（如显式屏蔽非法 token）可以在原则上完全消除部分类别错误，但会改变模型原始分布，本文将其留作后续工作。

\subsection{因果生成与 KV 缓存机制}

由于因果掩码确保历史隐藏状态在生成新 token 时不变，自回归推理可以缓存每层每个时间步的 K、V 矩阵，新增 token 只需与缓存做注意力点积，避免重复计算。KV cache 把单步推理复杂度从 $O(T^2)$ 降至 $O(T)$，在本文较长的序列设定下显著降低推理时间。LLaMA3 风格主干\cite{meta_2024_llama3} 在工程实现上也广泛采用 KV cache。

需要补充的是，KV cache 是一种纯粹的工程加速手段，它只消除了重复的注意力点积运算，并不改变长序列依赖本身的学习难度。本文 B-rep 序列中，某条边段的两个邻接面索引可能需要与几百 token 之前的对应面段保持一致，若序列化顺序完全随机，这种相关元素之间的实际距离会进一步拉大。第~3 章设计的拓扑感知排序就是希望在不改变模型结构的前提下，通过更合理的数据顺序把一部分几何—拓扑相关元素拉近，使自回归注意力更容易利用局部上下文。换言之，KV cache 解决“算得快”的问题，拓扑感知排序则尝试缓解“看得见”的问题，两者作用层面不同，本文同时使用。

\section{LLaMA3 风格模型组件}

LLaMA3\cite{meta_2024_llama3} 是 Meta 发布的开源大模型家族，其主干在 RMSNorm、SwiGLU、Rotary Position Embedding 与 Grouped Query Attention 等组件上沿用了近年来若干被反复验证的设计选择\cite{zhang_2019_rmsnorm,shazeer_2020_glu,su_2024_roformer,ainslie_2023_gqa}。本文采用 LLaMA3 风格的主干并不是为了套用大模型的容量，而是为了复用其中四个对生成稳定性、长序列适配和推理效率有直接帮助的组件——RMSNorm 简化了归一化统计、RoPE 让相对位置编码可以平滑外推、SwiGLU 提升了 FFN 在同参数预算下的表达能力、GQA 减小了 KV cache 的体积。下文按这四项依次给出本文使用的具体形式与超参数。

\subsection{RMSNorm}

RMSNorm 用每层激活的均方根代替 LayerNorm 中的均值—方差归一化：

\begin{equation}
    \text{RMSNorm}(x) = \frac{x}{\sqrt{\frac{1}{d}\sum_{i=1}^{d} x_i^2 + \epsilon}} \cdot \gamma,
    \label{eq:rmsnorm}
\end{equation}

其中 $\gamma$ 是可学习缩放向量。相较 LayerNorm，RMSNorm 去掉中心化操作，减少了一次按 batch 的均值统计，GPU 上数值更稳定，参数量也略小。本文主干每个 Transformer block 在注意力与 FFN 之前各使用一层 RMSNorm。

\subsection{旋转位置编码}

RoPE\cite{su_2024_roformer} 把位置信息编码为查询、键向量上的复数旋转：对相对位置 $m$，将向量的两个相邻通道 $(x_{2i}, x_{2i+1})$ 视为复数 $x_{2i} + j x_{2i+1}$，乘以单位复数 $e^{j m \theta_i}$，其中 $\theta_i = 10000^{-2i/d}$。RoPE 不显式插入位置嵌入向量，对长序列外推友好，并保持注意力点积的旋转不变性。本文主干仅对 Q、K 应用 RoPE，对 V 不作用。

\subsection{SwiGLU 门控前馈网络}

SwiGLU 用门控线性单元替换传统 FFN：

\begin{equation}
    \text{SwiGLU}(x) = \big(\text{Swish}(x W_1)\big) \odot (x W_2) \cdot W_3,
    \label{eq:swiglu}
\end{equation}

其中 $\text{Swish}(z) = z \cdot \sigma(z)$。在等参数预算下，SwiGLU 通常优于 GeLU 或 ReLU。本文主干每个 block 的 FFN 部分使用 SwiGLU。

\subsection{分组查询注意力机制}

GQA 让多个查询头共享同一组 K/V 头，介于多头注意力（MHA）与多查询注意力（MQA）之间。当 K/V 头组数等于查询头数时，结构退化为 MHA。当所有查询头共享同一组 K/V 时，则退化为 MQA。本文主干采用 8 个查询头和 4 个 K/V 头，每个 K/V 头由 2 个查询头共享，能在保持 KV cache 较小的同时维持注意力容量。

\section{CAD 生成质量评价指标}

\subsection{BRepCheck 几何有效性评价}

\begin{sloppypar}
OpenCASCADE 提供的 BRepCheck 校验器可以对输入 Shape 逐层检测面、边、顶点、Loop、Shell、Solid 的几何一致性与拓扑合法性。常见失败原因包括 Shell 含游离边、边的几何曲线不在面参数域上、取向不一致等。本文实现中调用该校验器作为统一有效性判据，当校验结果为有效时，生成样本才被认定为“有效”。
\end{sloppypar}

\subsection{生成结果三态统计口径}

本文主表口径与\autoref{tab:three-state-def} 一致，按“未产出 STEP / 产出但无效 / 产出且有效”三种状态对生成样本计数，并以“有效率（of saved）$= \mathrm{S3} / (\mathrm{S2} + \mathrm{S3})$”作为主要指标，即在所有成功写出 STEP 文件的样本中 BRepCheck 通过的比例，同时给出“有效率（of attempts）$= \mathrm{S3} / (\mathrm{S1} + \mathrm{S2} + \mathrm{S3})$”作为脚注口径。该口径与多数同类工作所报告的 OCC BRepCheck 通过率一致，便于与文献数据直接对照。

\subsection{评价指标适用性分析}

选择 BRepCheck 通过率作为主指标，是因为本文的核心目标不是三维外观展示，而是验证生成结果能否回到 CAD 内核可接受的 B-rep 实体。对于机械 CAD 场景，一个生成样本即使在 STL 网格中可视化正常，也可能因为 Shell 不闭合、边曲线不落在相邻面上或面取向冲突而无法参与后续布尔运算和尺寸编辑。BRepCheck 虽然不能覆盖所有工程可制造性要求，但它至少提供了一个比人工视觉判断更严格、也更容易复现实验的底线判据。

同时，BRepCheck 通过率也有其无法回答的几类问题，本文不应过度解读这一单一指标。最直接的局限是它只能回答“实体是否合法”，并不直接回答“实体是否美观、是否多样、是否接近训练集分布”——两个样本都能通过 BRepCheck，其中一个可能是简单长方体，另一个可能包含复杂孔洞和曲面过渡，二者在有效率统计中权重相同。其次，该指标对几何细节误差不敏感：只要误差未破坏内核合法性，就仍会被计为有效，因此即便实体在 0.5 mm 量级上偏离原始 CAD 几何，仍可能“通过”。再次，BRepCheck 的判定本身受重建容差、曲面拟合策略与 ShapeFix 配置影响，本文管线中 Sewing 容差和 ShapeFix 修复策略一旦改变，同一批生成样本得到的有效率会有所偏移，因此严格地说，该指标评估的不是孤立的神经网络性能，而是“模型输出 + 逆向重建管线”的整体闭环性能。本文报告有效率时会同时给出重建管线版本与采样配置，正是为了让该数字具备可追溯的口径。

基于上述原因，本文在第~5 章将 BRepCheck 通过率作为主指标，同时保留验证集交叉熵、PPL、S1/S2/S3 三态分布和序列长度分布等辅助证据。验证集交叉熵说明模型是否学到了训练分布中的 token 规律，三态分布说明失败发生在解析、重建还是有效性检查阶段，序列长度分布则说明当前上下文上限覆盖了怎样复杂度的样本。多指标共同使用，可以避免把单一有效率数字解释得过宽。
